% !PRINTED 557
\weberpage{557}

\S.\,154.

\begin{center}
\gesperrt{Die Gruppe der Idealclassen.}
\end{center}

Die Idealclassen können auf Grund des folgenden Satzes componirt werden:

1.~\gesperrt{Sind $\varphi$, $\psi$, $\varphi_1$, $\psi_1$ Functionale und $\varphi$ äquivalent mit $\varphi_1$, $\psi$ äquivalent mit $\psi_1$, so ist auch $\varphi\psi$ äquivalent mit $\varphi_1 \psi_1$.}

Die Richtigkeit hiervon ergiebt sich unmittelbar aus der Definition. Denn wenn $\varphi:\varphi_1$ und $\psi:\psi_1$ mit Zahlen associirt sind, so ist auch $\varphi\psi:\varphi_1\psi_1$ mit einer Zahl associirt.

Ebenso ergiebt sich auch der umgekehrte Satz:

2.~\gesperrt{Ist $\varphi$ äquivalent mit $\varphi_1$ und $\varphi\psi$ äquivalent mit $\varphi_1\psi_1$, so ist auch $\psi$ äquivalent mit $\psi_1$.}

Betrachten wir also zwei Idealclassen $A$, $B$, die auch identisch sein können, und bilden das Product $\varphi\psi$ irgend eines Functionals $\varphi$ aus $A$ und eines Functionals $\psi$ aus $B$, so ist die Classe $C$, in der das Product $\varphi\psi$ vorkommt, unabhängig von der Wahl von $\varphi$ und $\psi$, und die Classe $C$ ist durch die beiden Classen $A$, $B$ völlig bestimmt. \gesperrt{Wir nennen $C$ aus $A$ und $B$ componirt} und schreiben symbolisch

\weq{1}{C = AB = BA.}

Die Classe $C$ enthält alle Producte eines Elementes von $A$ mit einem Element von $B$, kann aber auch noch andere Functionale enthalten.

Da diese Composition aus der wahren Multiplication abgeleitet ist, so gelten auch die Gesetze der Multiplication für diese Composition, nämlich das commutative und das associative Gesetz.

Es folgt ferner aus dem Satze 2., dass, wenn $AB = AB_1$ ist, auch $B = B_1$ sein muss, und folglich erzeugen die Idealclassen bei dieser Composition eine \gesperrt{endliche Abel'sche Gruppe vom Grade $h$}, auf die wir alle Sätze anwenden können, die wir im zweiten Abschnitt dieses Bandes über solche Gruppen kennen gelernt haben.

Die Einheit dieser Gruppe ist die schon im \S.\,152 definirte Hauptclasse $O$, die ja, wie wir gesehen haben, den Repräsen-
%CONT>>
